\section{Simulación mediante autómatas celulares}

A lo largo del tiempo, se han propuesto múltiples variantes del modelo clásico de autómatas celulares, adaptadas a contextos no estructurados, probabilísticos, heterogéneos o continuos. A continuación, se describen algunas de estas extensiones para el modelado de sistemas complejos.

\begin{itemize}
	\item Los autómatas celulares en grafos (GCA, \textit{Graph Cellular Automata}) generalizan la estructura reticulada clásica permitiendo vecindades no uniformes: cada célula puede tener un número arbitrario y variable de vecinos. Esta flexibilidad los hace adecuados para modelar fenómenos sobre grafos, como redes sociales o flujos en infraestructura. Sin embargo, diseñar manualmente las reglas de transición en estos entornos es una tarea compleja, lo que ha motivado propuestas como los Graph Neural Cellular Automata (GNCA), que permiten aprender dichas reglas automáticamente \cite{neuripsGCA}.
	
	\item Los autómatas celulares probabilísticos (PCA, \textit{Probabilistic Cellular Automata}) extienden el modelo clásico al incorporar reglas de transición estocásticas. En lugar de una evolución determinista, los estados futuros dependen de distribuciones de probabilidad condicionadas por la vecindad, lo que permite capturar incertidumbre, ruido y comportamientos emergentes. Son especialmente útiles para simular dinámicas como propagación de epidemias o sistemas sociales, donde los mecanismos de transición no siempre están completamente definidos \cite{gacsPCA}.

	\item Los autómatas celulares no homogéneos (INCA, \textit{Inhomogeneous Cellular Automata}) permiten que diferentes células tengan diferentes funciones de transición. Para esto, se utilizan dos planos llamados \textit{data plane} el cual contiene la evolución del sistema; y \textit{control plane} el cual define las reglas de transición para cada célula. Esto permite estudiar el comportamiento y la interacción entre diferentes reglas que comparten espacio \cite{hartmanIHCA}.
	
	\item Los autómatas celulares continuos (CCA, \textit{Continuous Cellular Automata}) pueden ser continuos en los estados, en el tiempo y en el espacio \cite{wolfram2002}:
	\begin{itemize}
		\item Los estados continuos adquieren valores reales y no solo enteros.
		\item El tiempo continuo se define mediante una función que determina el intervalo de espera para actualizar la célula en función de su estado actual.
		\item El espacio continuo se obtiene al asignar diferentes pesos a las células dentro de la vecindad \cite{beckerCCA}.
	\end{itemize}
\end{itemize}

Si bien estas variantes han expandido el alcance de los CA, un desafío recurrente es la correcta definición de las reglas de transición para reproducir un fenómeno deseado. Este reto ha motivado la introducción de los Autómatas Celulares Neuronales (NCA, \textit{Neural Cellular Automata}), en los cuales la función de transición no se especifica explícitamente, sino que es aprendida mediante una Red Neuronal Artificial (ANN, \textit{Artificial Neural Network}). De este modo, las reglas locales emergen directamente de los datos, lo que facilita modelar dinámicas de alta complejidad.


Un primer avance en esta línea fue presentado por Gilpin \cite{gilpinNCA}, demostrando que una Red Neuronal Convolucional (CNN, \textit{Convolutional Neural Network}) puede aprender el comportamiento de un CA clásico. Su estudio destaca que tanto los CA como las CNN comparten dos propiedades fundamentales: (1) la actualización basada en una vecindad o \textit{kernel}, y (2) la aplicación simultánea de dicha función en todo el espacio.

En particular, Gilpin entreno una CNN con un \textit{kernel} de $3 \times 3$ (mismo tamaño que la vecindad de Moore) para imitar el comportamiento del \textit{glider}\footnote{Una de las configuraciones mas importantes por su capacidad de moverse en el espacio} en el Juego de la Vida (GoL, \textit{Game of Life}) de John Conway. Este trabajo estableció una conexión entre los CA y las CNN, mostrando que las ANN puede actuar como mecanismos de transición local aprendidos.

Posteriormente, esta línea de investigación fue extendida por Mordvintsev et al., quienes exploraron representaciones continuas para morfogénesis artificial (\textit{Growing NCA}), y por Niklasson et al., que aplicaron el enfoque a la generación de texturas dinámicas sobre superficies tridimensionales (\textit{Mesh NCA}). A su vez, otros autores han propuesto aplicaciones en robótica, clasificación de patrones y simulación de sistemas fotónicos, consolidando a los NCA como una herramienta para el modelado de sistemas autoorganizados.